\documentclass[12pt,a4paper]{article}
\usepackage{goyabase}

\usepackage{longtable}
\usepackage{calc}

% Pandoc compatibility
\providecommand{\tightlist}{%
  \setlength{\itemsep}{0pt}\setlength{\parskip}{0pt}}
\newcommand{\passthrough}[1]{#1}

\title{\textbf{Neo4j en Docker - UT8 NoSQL}}
\author{Departamento de Informática}
\date{\today}

\usepackage[spanish, es-noshorthands, es-tabla]{babel}
\usepackage[autostyle]{csquotes}
\begin{document}

\maketitle

\section{Neo4j en Docker - UT8 NoSQL}\label{neo4j-en-docker---ut8-nosql}

Este entorno permite trabajar con \textbf{Neo4j}, una base de datos de \textbf{Grafos} diseñada para gestionar relaciones complejas con un rendimiento inigualable.

\subsection{El Ecosistema de Datos de GloboTour: \enquote{SocialTour}}\label{ecosistema}
En GloboTour, usamos Neo4j para alimentar nuestra red social de viajeros, \textbf{SocialTour}. Aquí el valor no está en los datos aislados, sino en las conexiones entre ellos.

\begin{enumerate}
    \item \textbf{Nodos de Usuarios y Guías}: 
    \begin{itemize}
        \item \textit{Dato Real:} El nodo \texttt{(u:Usuario {nombre: 'Víctor', ciudad: 'Madrid'})} representa a un viajero.
    \end{itemize}
    \item \textbf{Relaciones Semánticas}: 
    \begin{itemize}
        \item \textit{Dato Real:} La relación \texttt{-[:SIGUE]->} conecta a dos usuarios.
        \item La relación \texttt{-[:VALORO {estrellas: 5}]->} conecta a un usuario con un guía turístico.
    \end{itemize}
    \item \textbf{El Poder del Grafo}:
    \begin{itemize}
        \item \textit{Dato Real:} Podemos responder en milisegundos a: \enquote{¿Qué guías han puntuado mis amigos que yo aún no conozco?}.
    \end{itemize}
\end{enumerate}

\subsection{Componentes}\label{componentes}

\begin{enumerate}
    \item \textbf{Neo4j Server}: Motor de grafos en el puerto \texttt{7687} (Bolt) y \texttt{7474} (HTTP).
    \item \textbf{Neo4j Browser}: Interfaz visual en el puerto \texttt{7474}.
\end{enumerate}

\subsection{Instrucciones de uso}\label{instrucciones-de-uso}

\subsubsection{Configuración del Entorno (Docker Compose)}\label{configuracion-del-entorno}

Este archivo define los servicios necesarios para la práctica:

\begin{lstlisting}[language=YAML]
services:
  neo4j:
    image: neo4j:latest
    container_name: neo4j-server
    ports:
      - "7474:7474"
      - "7687:7687"
    volumes:
      - neo4j_data:/data
      - neo4j_logs:/logs
      - ./src:/var/lib/neo4j/import
    environment:
      - NEO4J_AUTH=neo4j/password123
    restart: always

volumes:
  neo4j_data:
  neo4j_logs:
\end{lstlisting}

\subsubsection{Puesta en marcha}\label{puesta-en-marcha}

Para arrancar el servicio:

\begin{lstlisting}[language=bash]
docker compose up -d
\end{lstlisting}

\subsubsection{Cargar escenario por terminal (Recomendado)}\label{cargar-escenario-por-terminal-recomendado}

Para cargar el escenario de GloboTour directamente sin usar la web:

\begin{lstlisting}[language=bash]
cat src/setup_grafo.cypher | docker exec -i neo4j-server cypher-shell -u neo4j -p password123 --format plain 2>/dev/null
\end{lstlisting}

\subsubsection{Acceso a la Consola Interactiva (CLI)}\label{acceso-consola}

Si prefieres ejecutar comandos Cypher uno a uno directamente en la terminal (similar a MySQL):

\begin{lstlisting}[language=bash]
docker exec -it neo4j-server cypher-shell -u neo4j -p password123
\end{lstlisting}

\subsubsection{Cargar escenario por la Web}\label{cargar-escenario-por-la-web}

\begin{enumerate}
    \item Abre \url{http://localhost:7474}.
    \item Login: \texttt{neo4j} / \texttt{password123}.
    \item Copia el contenido de \texttt{src/setup\_grafo.cypher} en la barra superior y ejecuta.
\end{enumerate}

\subsection{Comandos sugeridos para probar}\label{comandos-sugeridos-para-probar}

\begin{itemize}
    \item \texttt{MATCH (n) RETURN n} (Ver todo el grafo).
    \item \texttt{MATCH (u:Usuario)-[r:AMIGO\_DE]->(f:Usuario) RETURN u,r,f} (Ver solo la red social).
\end{itemize}

\end{document}
